\documentclass[a4paper,11pt]{article}
\usepackage[utf8]{inputenc}\usepackage[T1]{fontenc}
\usepackage[ngerman]{babel}\usepackage{lmodern}\usepackage{microtype}
\usepackage{amsmath,amssymb}\usepackage{geometry}
\geometry{a4paper,left=2.2cm,right=2.2cm,top=2.0cm,bottom=2.0cm}
\usepackage{booktabs,array}\usepackage{xcolor}
\definecolor{bokug}{RGB}{0,135,60}\usepackage{tikz}
\usetikzlibrary{arrows.meta,calc}
\usepackage{fancyhdr}\pagestyle{fancy}
\renewcommand{\headrulewidth}{0.4pt}
\usepackage{enumitem}\setlength{\parindent}{0pt}\setlength{\parskip}{4pt}

\begin{document}
\fancyhead[L]{\small VU Mechanik – LAWI 100515 (PI)}\fancyhead[R]{\small SS 2026}\fancyfoot[C]{\thepage}
\begin{center}{\large\bfseries\color{bokug}Universität für Bodenkultur Wien}\\[2pt]{\large\bfseries VU Mechanik – LAWI 100515 (PI)}\\[4pt]Prüfung \textbf{Gruppe A} \hfill Datum: 26.06.2026\\[2pt]\rule{\linewidth}{0.4pt}\end{center}

\textbf{Gegebenes System}\par Das ebene System besteht aus Biegeträgern und einem Fachwerk, die durch Gelenke verbunden sind.
\begin{center}
\begin{tikzpicture}[scale=0.778,>=Stealth,line join=round]
  \coordinate (P0) at (0.000,0.000);
  \coordinate (P1) at (0.000,3.000);
  \coordinate (P2) at (0.000,6.000);
  \coordinate (TO1) at (-1.500,5.250);
  \coordinate (TU0) at (-1.500,6.750);
  \coordinate (TO2) at (-3.000,6.000);
  \draw[thick] (P2) -- (TO1);
  \draw[thick] (P2) -- (TU0);
  \draw[thick] (TO1) -- (TO2);
  \draw[thick] (TU0) -- (TO2);
  \draw[thick] (TO1) -- (TU0);
  \draw[line width=2.2pt] (P0) -- (P1);
  \draw[line width=2.2pt] (P1) -- (P2);
  \fill (TO1) circle (1.6pt);
  \fill (TU0) circle (1.6pt);
  \fill (TO2) circle (1.6pt);
  \draw[fill=white,line width=1pt] (P2) circle (3.2pt);
  \draw[line width=1.1pt] (-0.450,0.000) -- (0.450,0.000);
  \draw[line width=0.6pt] (-0.400,0.000) -- (-0.520,-0.160);
  \draw[line width=0.6pt] (-0.200,0.000) -- (-0.320,-0.160);
  \draw[line width=0.6pt] (0.000,0.000) -- (-0.120,-0.160);
  \draw[line width=0.6pt] (0.200,0.000) -- (0.080,-0.160);
  \draw[line width=0.6pt] (0.400,0.000) -- (0.280,-0.160);
  \draw[thick] (-3.000,6.000) -- (-2.740,6.360) -- (-3.260,6.360) -- cycle;
  \draw[thick] (-2.840,6.460) circle (0.06) (-3.160,6.460) circle (0.06);
  \draw[line width=1pt] (-2.580,6.540) -- (-3.420,6.540);
  \node[circle,fill=white,draw=gray!55,inner sep=0.5pt,font=\scriptsize] at (-0.750,5.625) {1};
  \node[circle,fill=white,draw=gray!55,inner sep=0.5pt,font=\scriptsize] at (-0.750,6.375) {2};
  \node[circle,fill=white,draw=gray!55,inner sep=0.5pt,font=\scriptsize] at (-2.250,5.625) {3};
  \node[circle,fill=white,draw=gray!55,inner sep=0.5pt,font=\scriptsize] at (-2.250,6.375) {4};
  \node[circle,fill=white,draw=gray!55,inner sep=0.5pt,font=\scriptsize] at (-1.500,6.000) {5};
  \draw[->,blue!70!black] (-0.620,3.000) -- (-0.060,3.000);
  \draw[->,blue!70!black] (-0.620,3.500) -- (-0.060,3.500);
  \draw[->,blue!70!black] (-0.620,4.000) -- (-0.060,4.000);
  \draw[->,blue!70!black] (-0.620,4.500) -- (-0.060,4.500);
  \draw[->,blue!70!black] (-0.620,5.000) -- (-0.060,5.000);
  \draw[->,blue!70!black] (-0.620,5.500) -- (-0.060,5.500);
  \draw[->,blue!70!black] (-0.620,6.000) -- (-0.060,6.000);
  \draw[blue!70!black,thick] (-0.620,3.000) -- (-0.620,6.000);
  \node[blue!70!black] at (-0.870,4.500) {$q$};
  \draw[->,red!80!black,line width=1.1pt] (-2.450,5.250) -- (-1.500,5.250);
  \node[red!80!black] at (-2.630,5.250) {$P$};
  \node[font=\small] at (0.000,0.000) [xshift=-7pt,yshift=7pt] {$A$};
  \node[font=\small] at (0.000,3.000) [xshift=-7pt,yshift=7pt] {$D$};
  \node[font=\small] at (0.000,6.000) [xshift=-7pt,yshift=7pt] {$C$};
  \node[font=\small] at (-3.000,6.000) [xshift=-7pt,yshift=7pt] {$B$};
\end{tikzpicture}
\end{center}
\textbf{Gegebene Größen:}\quad $a=3$\,m\quad $\ell=1,5$\,m\quad $h=1,5$\,m\quad $q=5$\,kN/m\quad $P=15$\,kN
\par\medskip
\textbf{Aufgabenstellung}\par
\begin{enumerate}[label=\textbf{\arabic*.}]
  \item Bestimmen Sie den Grad der statischen Bestimmtheit.
  \item Berechnen Sie alle Auflagerreaktionen und Gelenkkräfte.
  \item Ermitteln Sie alle Stabkräfte des Fachwerks (Zug/Druck).
  \item Zeichnen Sie die Schnittgrößen $N$, $V$, $M$ der Biegeträger.
\end{enumerate}
\end{document}